\section{硬件微观机制深讲：状态机、缓存一致性和 I/O 事务}

上一节把整机硬件铺开，本节继续往下钻。学习操作系统最容易卡住的地方，是知道“有 cache、有中断、有 DMA”，但不知道这些机制内部怎样改变状态，为什么内核要按某种顺序写代码。本节按实现粒度讲：CPU 如何提交指令，cache line 如何在核心之间转移，TLB 如何失效，PCIe/NVMe 如何完成一次事务，IOMMU 如何把设备也放进地址空间。

\subsection{同步时序：为什么硬件要靠时钟推进}

组合逻辑只把输入变成输出，不能记住过去。寄存器在时钟边沿采样输入，构成下一周期的状态。CPU、内存控制器、DMA 引擎、NVMe 控制器、网卡队列，本质上都是许多状态机协同推进。

\begin{lstlisting}
state machine template:
  current_state is stored in registers
  combinational logic computes next_state from inputs
  on clock edge:
    if reset:
      current_state = RESET_STATE
    else:
      current_state = next_state
\end{lstlisting}

OS 里的“状态不可跳跃”来自这里。设备 reset 后必须等 ready bit，是因为设备内部状态机还没从 RESET 走到 READY；写 doorbell 后要等 completion，是因为设备只记录了“有新请求”，数据搬运、错误检查和完成写回还要经历多个状态；中断处理必须给 EOI，是因为中断控制器也有 pending、in-service、masked 等状态。

\subsection{流水线不是五个名词，而是一组提交规则}

教科书常把流水线写成取指、译码、执行、访存、写回。真实 CPU 可能有更复杂的前端、重命名、发射、执行、回写、提交，但对 OS 最关键的是提交规则：架构状态必须按程序顺序提交，异常必须精确。

\begin{lstlisting}
out-of-order core sketch:
  fetch instruction bytes
  decode into micro-ops
  rename architectural registers to physical registers
  dispatch micro-ops into queues
  execute ready micro-ops out of order
  write results to physical registers
  retire completed instructions in program order
\end{lstlisting}

“乱序执行”不意味着用户程序看见乱序结果。CPU 可以内部先执行后面的 load 或 ALU 操作，但提交到架构寄存器、异常点和内存可见语义时必须遵守 ISA 规定。缺页异常能返回原指令重试，依赖的就是精确异常：出错指令之前的指令已经提交，之后的指令没有提交。

\subsection{精确异常的例子：load 缺页}

假设指令流中第三条 load 缺页，第四条加法可能已经在乱序引擎里算过，但不能对 OS 可见。

\begin{lstlisting}
inst1: add r1, r2, r3       committed
inst2: store [x], r1        committed or ordered by memory rules
inst3: load r4, [bad_va]    raises page fault
inst4: add r5, r6, r7       may have executed speculatively

architectural view at exception:
  saved RIP points to inst3
  inst1 and inst2 are visible as completed
  inst3 has not produced final architectural result
  inst4 is not architecturally visible
\end{lstlisting}

内核处理缺页时才能做出可靠判断：若是合法懒分配，就建立 PTE 后返回 inst3 重试；若是非法访问，就向进程发送信号。没有这个契约，page fault 不能成为虚拟内存的基础。

\subsection{寄存器重命名和上下文切换}

CPU 内部可能有比 ISA 更多的物理寄存器。寄存器重命名把软件可见寄存器映射到物理寄存器，消除很多假依赖。上下文切换时，OS 不需要保存所有物理寄存器，只需要保存 ISA 规定的软件可见状态。

\begin{longtable}{p{0.24\linewidth}p{0.34\linewidth}p{0.3\linewidth}}
\toprule
状态 & 谁管理 & 上下文切换关系 \\
\midrule
架构通用寄存器 & OS 必须保存恢复 & 用户程序可见，必须准确 \\
浮点/SIMD 状态 & OS 按 ABI 和策略保存 & 状态大，可能惰性或按需保存 \\
物理寄存器池 & CPU 微架构管理 & OS 不直接保存 \\
分支预测器状态 & CPU 微架构管理 & 性能和侧信道相关 \\
cache/TLB & CPU 和 OS 协作维护 & 不保存，但切换影响命中率 \\
\bottomrule
\end{longtable}

这解释了为什么上下文切换不仅是保存几十个寄存器。它还会改变 cache/TLB 热状态、分支预测历史、内存局部性和 NUMA 访问路径。调度器选择任务时，实际在权衡“马上公平”和“保留热状态”。

\subsection{cache line 状态：以 MESI 思想理解一致性}

多核系统中，每个 CPU 有自己的 cache。同一个物理 cache line 可能同时存在于多个核心。硬件一致性协议维护每条 line 的状态。不同 CPU 具体协议不同，但用 MESI 思想足以理解 OS 成本。

\begin{longtable}{p{0.16\linewidth}p{0.36\linewidth}p{0.36\linewidth}}
\toprule
状态 & 含义 & 典型转移 \\
\midrule
Modified & 本核有最新数据，内存可能旧 & 其他核读时必须提供或写回 \\
Exclusive & 本核独占且干净 & 本核写入后变 Modified \\
Shared & 多个核可读共享 & 某核写入前要让其他副本失效 \\
Invalid & 本核没有有效副本 & 读写时要重新获取 line \\
\bottomrule
\end{longtable}

两个 CPU 写同一 cache line 上的不同变量时，line 会在两个核心间来回转移。程序逻辑上写的是两个变量，硬件上争的是同一 line 的独占权。这就是伪共享。内核的 per-CPU 数据、cacheline 对齐、分片计数器，本质都是减少 line 状态来回迁移。

\subsection{锁的真实成本：cache line 所有权迁移}

自旋锁通常是一个小整数。获取锁时，CPU 执行原子读改写，必须拿到锁变量所在 cache line 的独占权。很多 CPU 同时争锁时，这条 line 在核心之间移动，互联流量和延迟迅速上升。

\begin{lstlisting}
test_and_set_lock(lock):
  while atomic_exchange(lock, 1) == 1:
    cpu_relax()

unlock(lock):
  store_release(lock, 0)
\end{lstlisting}

这个简单锁在低竞争时很好，高竞争时很差。排队自旋锁把“所有 CPU 抢同一变量”变成“每个 CPU 等自己的节点”，减少缓存抖动。

\begin{lstlisting}
queued_spin_lock:
  node = per_cpu_node()
  pred = atomic_xchg(lock.tail, node)
  if pred != null:
    node.locked = true
    pred.next = node
    while node.locked:
      cpu_relax()
  // enter critical section
\end{lstlisting}

Linux 的 queued spinlock 和 MCS 思想相关。这里的重点不是记算法名，而是理解：高竞争锁慢的主要原因常常不是临界区代码，而是 cache line 所有权在 CPU 间转移。

\subsection{内存顺序：store buffer 为什么会改变可见顺序}

CPU 为了性能会把 store 先放进 store buffer，稍后再对其他核心可见。单线程看不出问题，因为本 CPU 之后的 load 可以从 store buffer 转发；多核就可能看到不同顺序。

\begin{lstlisting}
CPU0:
  data = 42
  flag = 1

CPU1:
  if flag == 1:
    assert(data == 42)
\end{lstlisting}

在弱顺序模型或编译器重排下，CPU1 看到 \code{flag == 1} 时不一定能看到 \code{data == 42}，除非使用 release/acquire 或锁。正确发布对象要写成：

\begin{lstlisting}
CPU0:
  data = 42
  store_release(flag, 1)

CPU1:
  if load_acquire(flag) == 1:
    assert(data == 42)
\end{lstlisting}

OS 内核里同样适用。唤醒等待者前要先改变等待条件；发布文件对象前要先初始化字段；设备 doorbell 前要先让描述符对设备可见。屏障的含义永远要和“谁观察什么”一起写。

\subsection{TLB 和 page walk cache}

TLB 缓存虚拟页到物理页的翻译。现代 CPU 还可能有 page walk cache，缓存页表中间层结果。修改页表后，内存里的 PTE 变了，不代表 CPU 立即不用旧翻译。OS 必须按架构要求失效相关 TLB。

\begin{lstlisting}
page walk for a 4-level page table:
  read top-level entry from page table root
  read next-level entry
  read next-level entry
  read final PTE
  check present and permissions
  insert translation into TLB
\end{lstlisting}

大页减少页表层级和 TLB 项数量，但权限粒度变粗。一个 2MiB 大页内只要有 4KiB 要改权限，就可能拆成 512 个小页。拆页会造成页表分配、TLB flush 和锁竞争，所以透明大页可能提升吞吐，也可能造成延迟尖峰。

\subsection{TLB shootdown 的正确性边界}

释放物理页前必须保证没有 CPU 还能通过旧 TLB 访问它。否则这页可能已经分配给别的进程或内核对象，旧翻译会变成越权访问。

\begin{lstlisting}
safe unmap:
  remove PTE or clear permission
  flush local TLB entry
  send shootdown IPI to CPUs that may use mm
  remote CPUs invalidate entry and acknowledge
  wait for acknowledgements
  now old physical page may be freed or reused
\end{lstlisting}

这里的“may use mm”不是随便猜。内核要跟踪哪些 CPU 正在或最近运行某个地址空间，避免给所有 CPU 发 IPI。过度 shootdown 会让大机器扩展性下降；少发一个 shootdown 则是安全 bug。

\subsection{中断控制器状态：pending、masked、in-service}

设备触发中断后，中断控制器通常不是直接丢给 CPU 就结束。它会记录 pending，检查 mask 和优先级，把 vector 送到目标 CPU。CPU 进入 handler 后，中断可能处于 in-service 状态；handler 结束时需要 EOI 或等价操作。

\begin{lstlisting}
interrupt controller:
  if device asserts interrupt:
    mark vector pending
  if not masked and priority allows:
    deliver vector to target CPU
    mark vector in_service
  after handler sends EOI:
    clear in_service
    deliver next pending interrupt if any
\end{lstlisting}

若驱动没有清设备状态就 EOI，中断可能立刻再次触发，形成中断风暴。若 handler 清了设备状态但忘了 EOI，后续同类中断可能被阻塞。若 mask 过久，completion 延迟会飙升。中断处理顺序必须同时满足设备状态机和中断控制器状态机。

\subsection{MSI-X：中断也是一次内存写消息}

传统线中断像一根线被拉高。MSI/MSI-X 更像设备向特定地址写入一条消息，中断控制器把这条消息解释成某个 vector。这样设备可以有多个 vector，适合多队列网卡和 NVMe。

\begin{lstlisting}
MSI-X setup:
  kernel allocates interrupt vectors
  kernel programs device MSI-X table
  each queue gets vector and CPU affinity
  device writes MSI-X message when queue completes
  target CPU enters corresponding handler
\end{lstlisting}

多队列设备性能依赖这个映射。若 64 个 RX/TX 队列都映射到一个 CPU，硬件有多队列也发挥不出来；若队列、CPU 和 NUMA 不匹配，DMA 数据会跨节点移动。驱动初始化时的 vector 分配不是细枝末节，而是性能结构。

\subsection{PCIe 事务：BAR、doorbell 和 completion}

PCIe 设备通过 BAR 暴露 MMIO 寄存器。CPU 写 doorbell 时，写请求会经过 root complex、switch、链路，到设备。写可能是 posted，CPU 不等待设备真正处理完成。设备完成 DMA 后，再通过 completion queue 和 MSI-X 通知 CPU。

\begin{lstlisting}
CPU submit to PCIe device:
  CPU writes descriptors in RAM
  dma_wmb()
  CPU MMIO-writes doorbell BAR
  posted write travels through PCIe fabric
  device receives doorbell
  device DMA-reads descriptors
  device performs work
  device DMA-writes completion
  device sends MSI-X message
\end{lstlisting}

这个路径解释了为什么驱动不能把 MMIO 写当函数调用。doorbell 返回不代表设备完成；甚至不一定代表设备已经处理。若驱动需要确认某个 posted write 到达，常见做法是在平台规则允许时读回一个寄存器形成排序点。

\subsection{NVMe 队列内部：phase tag 和 head/tail}

NVMe submission queue 由主机写、设备读；completion queue 由设备写、主机读。completion queue 常用 phase tag 区分“这一格是新一轮写入还是上一轮旧值”。队列满、队列空、环绕时都要靠 head/tail 和 phase 正确判断。

\begin{lstlisting}
completion poll:
  cqe = cq[head]
  if cqe.phase != expected_phase:
    return no_new_completion
  process cqe
  head = head + 1
  if head == queue_size:
    head = 0
    expected_phase = !expected_phase
  write new head doorbell
\end{lstlisting}

如果没有 phase tag，环形队列回绕后，软件无法可靠区分旧 completion 和新 completion。这个细节体现了硬件队列的共同问题：数组槽位会复用，所以必须有 generation、phase、sequence 或所有权位。

\subsection{IOMMU 页表和 IOTLB}

IOMMU 给设备也建立翻译缓存。设备发出 IOVA，IOMMU 查页表得到物理地址，并检查读写权限。为了速度，IOMMU 也会有 IOTLB。撤销 DMA 映射后，必要时也要失效 IOTLB，否则设备可能继续用旧映射。

\begin{lstlisting}
dma_unmap:
  remove IOVA -> physical mapping from IOMMU page table
  invalidate IOTLB for this device/domain/range
  only after invalidation may page be returned to allocator
\end{lstlisting}

这和 CPU TLB shootdown 是同一类正确性问题，只是访问者从 CPU 变成设备。设备侧隔离必须考虑 IOMMU group：若多个设备在硬件上不能区分请求者 ID，就不能把它们当成完全隔离的安全主体。

\subsection{非一致 DMA 平台上的 cache 维护}

有些平台上，设备 DMA 与 CPU cache 不自动一致。CPU 写给设备读的 buffer 前，需要 clean cache，把 dirty line 推到设备可见层；设备写给 CPU 读的 buffer 后，需要 invalidate cache，避免 CPU 读旧 line。

\begin{lstlisting}
TO_DEVICE:
  CPU writes buffer
  clean cache lines for buffer
  device DMA reads memory

FROM_DEVICE:
  device DMA writes memory
  invalidate CPU cache lines for buffer
  CPU reads fresh data
\end{lstlisting}

这就是为什么驱动必须使用 DMA API，而不是直接把物理地址写进描述符。即使本书主线是 x86-64，DMA API 仍然负责表达方向、生命周期、IOMMU 映射和必要屏障。直接绕开 API 的驱动，可能在一个测试机上偶然工作，却在启用 IOMMU、不同 cache 属性或不同设备 DMA mask 时随机损坏数据。

\subsection{存储介质内部状态：FTL 的写放大}

SSD 对外暴露逻辑块，对内管理 NAND 页和擦除块。覆盖写不能直接改旧页，通常写到新物理页，再把旧页标为无效。垃圾回收把仍有效的页搬到新块，然后擦除旧块。主机随机写越多，内部搬迁越多，写放大越大。

\begin{lstlisting}
FTL overwrite logical block L:
  old_phys = map[L]
  new_phys = allocate_free_flash_page()
  program new_phys with data
  map[L] = new_phys
  mark old_phys invalid
  later garbage_collect blocks with many invalid pages
\end{lstlisting}

TRIM/discard 告诉 SSD 某些逻辑块不再包含有效文件数据，FTL 可以把对应物理页当作无效页，降低后续垃圾回收成本。文件系统、块层和设备固件因此共同决定长期写入性能。

\subsection{可靠提交：cache flush、FUA 和日志顺序}

存储设备可能有易失缓存。块层写完成不一定表示掉电后存在。可靠提交要约束顺序：先写日志或数据，再 flush 或使用 FUA，最后提交元数据。设备必须向上层报告它支持哪些 flush 语义。

\begin{lstlisting}
journal transaction:
  write descriptor block
  write data or metadata blocks
  issue flush if device cache is volatile
  write commit block with FUA or followed by flush
  transaction is durable after commit reaches stable media
\end{lstlisting}

若 flush 被设备虚假完成，文件系统再正确也无法保证崩溃恢复。若文件系统漏掉目录 fsync，应用的原子替换协议也会有洞。可靠性是应用、文件系统、块层和设备共同实现的契约。

\subsection{硬件算法小结}

\begin{rubricbox}
\begin{itemize}
\tightlist
\item 能解释精确异常为什么让缺页可恢复。
\item 能用 MESI 思想解释伪共享和锁竞争。
\item 能说明 release/acquire 约束的是跨 CPU 可见顺序。
\item 能写出 TLB shootdown 和 IOMMU unmap 的安全顺序。
\item 能解释 MSI-X、多队列和 IRQ affinity 为什么影响 NVMe/网卡性能。
\item 能说明 NVMe completion queue 为什么需要 phase tag。
\item 能解释非一致 DMA 平台为什么必须 clean/invalidate cache。
\item 能把 SSD FTL、TRIM、flush 和文件系统日志顺序联系起来。
\end{itemize}
\end{rubricbox}
